%======================================================================
\section{Stackelberg 共识博弈}
\label{sec:game}
%======================================================================

\subsection{博弈公式化}

我们将多智能体共识公式化为一个 Stackelberg 博弈 $\mathcal{G} = (\mathcal{L}, \mathcal{F}, \Theta, \mathcal{S}, U_L, U_F)$。

\paragraph{参与者。} \emph{领导者}（leader）$\mathcal{L}$ 是机制设计者（双层优化器）。\emph{追随者}（followers）$\mathcal{F} = \{1, \ldots, n\}$ 是大语言模型（LLM）智能体。

\paragraph{策略空间。} 领导者选择机制参数 $\theta \in \Theta \subset \mathbb{R}^{21}$，控制 AEIC 质量信号如何映射到共识权重，以及如何调节集体行为动态。每个追随者 $i$ 选择一个策略 $s_i \in \{s^{\text{ind}}, s^{\text{syc}}\}$：
\begin{itemize}
    \item $s^{\text{ind}}$（独立推理）：以成本 $C_{\text{ind}}$ 产生高 IC、EC、EQ；
    \item $s^{\text{syc}}$（阿谀奉承）：产生较低的 IC/EC，但可能以较低的成本 $C_{\text{syc}} < C_{\text{ind}}$ 获得高 EV。
\end{itemize}

\paragraph{时间线。}
\begin{enumerate}
    \item 领导者承诺参数 $\theta$（机制规则）。
    \item 每个追随者观察到 $\theta$ 并选择策略 $s_i$。
    \item 智能体生成 AEIC 输出；计算质量信号 $\tau_i$。
    \item 机制通过评分规则计算权重 $w_i(\tau_i; \theta)$。
    \item 加权投票产生共识决策。
\end{enumerate}

\paragraph{效用函数。} 追随者的效用为：
\begin{equation}
    U_i(s_i; \theta) = w_i(\tau_i(s_i); \theta) \cdot R - C(s_i)
    \label{eq:follower_utility}
\end{equation}
其中 $R > 0$ 是被采纳的奖励。领导者最大化：
\begin{equation}
    \max_{\theta \in \Theta} \; \text{Accuracy}(\theta) + \lambda \cdot \text{IC}(\theta)
    \label{eq:leader_obj}
\end{equation}

\subsection{评分规则}
\label{sec:scoring}

21 个参数 $\theta$ 控制三个方面：

\begin{itemize}
    \item \textbf{出价构建}（13个参数）：IC、EC、EQ、EV 信号如何组合成智能体“出价”（对齐意愿 $\alpha_i$、修正成本、历史对齐度 $h_i$）。
    \item \textbf{Stackelberg 动态}（4个参数）：共识能量系数 $\lambda_c$、预算惩罚 $\mu$、学习率 $\eta$ 和同伴压力 $\gamma$。
    \item \textbf{动态调节}（4个参数）：阿谀奉承检测阈值 $\tau_h$、僵局阈值 $\tau_l$、EV 衰减率 $\alpha$ 和压力调整率 $\beta$（见第~\ref{sec:dynamic} 节）。
\end{itemize}

具体而言，对于智能体 $i$：
\begin{align}
    \alpha_i &= \max\!\big(\theta_{\text{floor}}, \; \text{IC}_i^{\theta_{\text{ic}}} \cdot \text{EQ}_i^{\theta_{\text{eq}}} \cdot (\theta_{\text{ev\_base}} + \theta_{\text{ev\_scale}} \cdot \text{EV}_i)\big) \\
    h_i &= \theta_{h_1} \cdot \text{IC}_i + \theta_{h_2} \cdot \text{EC}_i + \theta_{h_3} \cdot \text{EQ}_i + \theta_{h_4} \cdot \text{EV}_i
\end{align}

这些出价通过一个分布式共识博弈进行处理，该博弈通过由通信拓扑定义的图拉普拉斯矩阵引入同伴压力，并经过迭代更新后产生最终权重 $w_i(\theta)$。

\subsection{激励兼容性指标}
\label{sec:ic_metric}

我们通过组间权重比较来定义激励兼容性：

\begin{definition}[激励兼容性指标]
对于每个问题 $q$，按质量得分 $\phi(\tau_i) = (\text{IC}_i + \text{EC}_i + \text{EQ}_i)/3$ 对智能体进行排序。设 $H_q$ 为排名前 2 的智能体，$L_q$ 为排名后 2 的智能体。定义：
\begin{equation}
    \text{IC}(\theta) = \frac{1}{|Q|} \sum_{q \in Q} \mathbb{1}\!\left[\frac{1}{|H_q|}\sum_{i \in H_q} w_i(\theta) > \frac{1}{|L_q|}\sum_{j \in L_q} w_j(\theta)\right]
    \label{eq:ic_metric}
\end{equation}
\end{definition}

该指标衡量了较高质量的推理获得更高共识权重的问题比例，捕捉了激励兼容性的本质。

\subsection{双层优化}
\label{sec:bilevel}

Stackelberg 博弈通过双层优化求解：
\begin{align}
    \textbf{上层 (领导者):} \quad & \max_{\theta \in \Theta} \; \text{Accuracy}(\theta) + \lambda \cdot \text{IC}(\theta) \label{eq:upper} \\
    \textbf{下层 (追随者):} \quad & w_i = f(\tau_i; \theta) \quad \forall i \label{eq:lower}
\end{align}

给定 $\theta$，下层问题是确定性的（智能体的 AEIC 输出是离线固定的），从而简化为单层优化，这可以通过差分进化算法~\cite{storn1997differential} 来实现。针对每个拓扑优化的 $\theta^*_T$ 使用种群大小 21（参数数量的 1 倍）和 30 代进行优化，然后通过 5 折分层交叉验证进行评估，以防止训练/测试数据泄露。我们注意到，Stackelberg 机制提供了形式上的 IC 保证（定理~\ref{thm:ic}），但实验（第~\ref{sec:exp} 节）表明，在同质 LLM 智能体下，质量加权投票的性能并未超越多数投票——瓶颈在于质量信号有限的判别能力，而不是机制设计。

\subsection{在线 Stackelberg 扩展}
\label{sec:online_sg}

在多轮拓扑（辩论、分层）中，机制\emph{在线}（online）运行：在第 $t$ 轮和第 $t{+}1$ 轮之间，它从第 $t$ 轮的输出计算中间 AEIC 信号，运行 Stackelberg 博弈以获得当前权重 $w_i^{(t)}$，并将质量反馈注入到智能体提示（prompt）中：
\begin{equation}
    \text{context}_i^{(t+1)} = \text{context}_i^{(t+1)} \cup \{(\text{IC}_i^{(t)}, \text{EC}_i^{(t)}, w_i^{(t)})\}
\end{equation}
这将博弈从事后重加权转变为实时行为塑造：智能体观察到它们当前的质量评估，并可以据此调整推理投入。星形拓扑本质上是在线的——通过 Stackelberg 计算的叶节点信任分数被注入到中心节点的综合上下文中。

结合第~\ref{sec:bilevel} 节的静态优化，这建立了一个两层 Stackelberg 演化机制：静态优化发现要抑制\emph{什么}行为（例如，在信息共享拓扑上对 EV 赋予接近零的权重），而在线反馈传达了每个智能体当前表现\emph{如何}。然而，在线扩展在不考虑集体行为模式的情况下注入相同的质量反馈——它无法区分阿谀奉承的从众行为与富有成效的收敛。第~\ref{sec:dynamic} 节引入了第三层来解决这个限制。

\subsection{动态阿谀奉承调节}
\label{sec:dynamic}

在线 Stackelberg 反馈（第~\ref{sec:online_sg} 节）在轮次之间注入质量信号，但不论集体行为如何均应用统一的干预。在实践中，多轮协商展现出明显的病态模式：\emph{阿谀奉承式收敛}（智能体迅速放弃独立立场，追随正在形成的多数派）和\emph{僵局}（智能体固执地维持初始立场，而不与对立证据进行交锋）。我们引入了一个动态调节层，用于监控集体立场的动态变化并应用差异化的干预措施。

\paragraph{行为指标。}
在第 $t{-}1$ 轮和第 $t$ 轮之间，我们从智能体的结论 $c_i^{(t)}$（可靠地从 AEIC 的结论字段中提取）计算三个指标：
\begin{align}
    \text{SCR}^{(t)} &= \frac{1}{n}\sum_{i=1}^{n} \mathbb{1}\!\left[c_i^{(t)} \neq c_i^{(t-1)}\right] \label{eq:scr} \\
    \text{MR}^{(t)} &= \max\!\left(\frac{|\{i : c_i^{(t)} = \text{yes}\}|}{n},\; \frac{|\{i : c_i^{(t)} = \text{no}\}|}{n}\right) \label{eq:mr} \\
    \text{CS}^{(t)} &= \left|\text{MR}^{(t)} - \text{MR}^{(t-1)}\right| \label{eq:cs}
\end{align}
其中，SCR（立场改变率，Stance Change Rate）衡量改变了结论的智能体比例，MR（多数比例，Majority Ratio）衡量主导立场的支配程度，CS（收敛速度，Convergence Speed）捕捉多数派增长的速度。

\paragraph{状态分类。}
系统将每次轮间过渡分类为以下三种状态之一：
\begin{itemize}
    \item \textbf{阿谀奉承}（$\text{SCR} \geq \tau_h$）：很大比例的智能体改变了立场，表明潜在的从众行为，即智能体放弃独立推理以追随正在形成的多数派。
    \item \textbf{僵局}（$\text{SCR} \leq \tau_l \wedge \text{MR} < 0.8$）：很少有智能体改变立场，且不存在明显的多数派，表明阵营分明，缺乏富有成效的证据交流。
    \item \textbf{健康}（其他）：立场呈适度演变，出现逐渐收敛或稳定的分歧——不需要干预。
\end{itemize}
阈值 $\tau_h$、$\tau_l$ 以及下文的响应参数 $\alpha$、$\beta$ 包含在可学习参数向量 $\theta$ 中，将基础的 17 个参数扩展到 21 维。双层优化（第~\ref{sec:bilevel} 节）与基础 Stackelberg 参数一起联合发现最佳检测阈值和响应强度。

\paragraph{自适应响应。}
在检测到非健康状态时，机制通过两个渠道做出响应：
\begin{enumerate}
    \item \textbf{参数调整}：对于阿谀奉承状态，证据共识权重发生衰减：$\theta_{\text{ev\_scale}} \leftarrow \theta_{\text{ev\_scale}} \cdot (1 - \alpha)$，从而抑制这种奖励从众的信号。同伴压力和共识能量同样降低：$\gamma \leftarrow \gamma \cdot (1 - \beta)$，$\lambda_c \leftarrow \lambda_c \cdot (1 - \beta)$。对于僵局，这些参数则被\emph{增加}以鼓励互动。$\alpha \in [0.1, 0.9]$ 和 $\beta \in [0.1, 0.6]$ 均由领导者进行优化。
    \item \textbf{差异化提示}：下一轮的上下文中附加了特定状态的指令：反阿谀奉承提示强调独立推理而非随波逐流；交流提示鼓励与相反的证据进行辩论；中性提示提供标准的协商指导。
\end{enumerate}

这就完成了三层 Stackelberg 演化机制：\emph{静态}优化发现需要抑制哪些行为（第~\ref{sec:bilevel} 节），\emph{在线}反馈传达个人质量评估（第~\ref{sec:online_sg} 节），\emph{动态}调节实时检测和纠正集体行为病态。核心见解是 Stackelberg 博弈的价值在于\emph{集体行为调节}——塑造群体的协商方式——而不是基于质量选择个人贡献。
